\documentclass[10pt,a4paper]{amsart}
\usepackage[margin=19mm]{geometry}
\usepackage{amsmath,amssymb,mathtools}
\usepackage[scr=rsfs,cal=euler]{mathalfa}
\usepackage{xfrac}
\usepackage[unicode,hidelinks,bookmarks=false]{hyperref}
\newcommand{\ie}{i.e.\ }
\newcommand{\cf}{cf.\ }
\newcommand{\resp}{respectively\ }
\newcommand{\Subsection}{\subsection}
\providecommand{\nobreakdash}{\nobreak\hbox{-}}
\setlength{\emergencystretch}{2em}
\multlinegap0pt
\usepackage{bm}
\usepackage{amssymb}
\usepackage{xcolor}
\newtheorem{corollary}{Corollary}
\title{Dimension Reduction and Asymptotic Approximation of Reactive Transport in a Bulk Domain with a Branched Thin Fracture}
\author{Taras Mel'nyk$^\flat$ \ \& \ Christian Rohde$^\natural$}
\date{Source: arXiv:2607.29605v1 (CC BY 4.0)}
\begin{document}
\maketitle

\noindent\emph{Source and licence.} Dimension Reduction and Asymptotic Approximation of Reactive Transport in a Bulk Domain with a Branched Thin Fracture. Version arXiv:2607.29605v1 (CC BY 4.0), distributed by arXiv under the Creative Commons Attribution 4.0 International licence, http://creativecommons.org/licenses/by/4.0/, as recorded in the arXiv licence metadata for this exact version and shown on the abstract page https://arxiv.org/abs/2607.29605v1. Full licence notice: \texttt{licenses/arxiv-2607.29605-CC-BY-4.0.txt}.\par\smallskip

\noindent\textit{Editorial scope.} Counted unit: the article's corollary corollarity-2-3 (source0.tex lines 2144--2152). The article's complete original proof of it is delivered with it (source0.tex lines 2153--2207, one \texttt{proof} environment of 55 lines). No mathematics is written, repaired or re-derived: the statement and the proof are the article's own lines, in the article's own notation, and no retained line is substituted or re-flowed.  Retained uncounted as the source-local context the proof needs: lines 2117--2122; lines 2128--2133.  Nothing was omitted from any retained span, so the package carries no omission record.  No retained line contains a citation, so the wrapper carries no bibliography and no bibliography entry.  No retained line contains a figure environment, an image-inclusion command or drawing code, so no image is delivered and the package declares no asset dependency.  Editorial formatting only, in addition: an AMS \texttt{amsart} preamble that restates the article's own declarations and macros used by the retained lines, namely the environments \texttt{corollary} on separate counters, together with the standard packages those lines need.  The retained environments print in this document as Corollary 1 in order of appearance, whereas the article prints the same items with the numbering of its own full document.  The complete native source of the article as distributed in the arXiv e-print for this exact version is bundled unchanged as \texttt{sources/206503/source0.tex}.
		\begin{equation}\label{asym-5}
			u^{(i)}_1(x,t)
			= C^{(i)}_1(t)\,\log r + \widetilde{C}^{(i)}_1(t)
			+ \mathcal{O}\!\left(r^{\delta_i}\right)
			\quad \text{as } \ \  r=\sqrt{x_1^2+x_2^2}\to 0,
		\end{equation}

		\begin{equation}\label{asym-deriv-r}
			\partial_r u^{(i)}_1(x,t)
			= \frac{C^{(i)}_1(t)}{r}
			+ \mathcal{O}\!\left(r^{\delta_i-1}\right)
			\quad \text{as } \ \ r\to 0,
		\end{equation}

	\begin{corollary}\label{corollarity-2-3}
		The normal derivative of $u^{(i)}_1$ to the sides $\Gamma_\varepsilon^{(i,i)}$ and $\Gamma_\varepsilon^{(i,j)}$ satisfies
		\begin{equation}\label{asym-normal-line-final}
			\partial_{\boldsymbol{\nu}_\varepsilon} u^{(i)}_1(x,t)
			= C^{(i)}_1(t)
			+ \mathcal{O}\!\left(\varepsilon^{\delta_i/2}\right)
			\quad \text{as } \ \ r_\varepsilon = |x| = \mathcal{O}\!\left(\varepsilon^{1/2}\right) \ \ \text{and} \ \ \varepsilon\to 0.
		\end{equation}
	\end{corollary}

	\begin{proof}
		We restrict ourselves to the case $i = 1$ and the side $\Gamma_\varepsilon^{(1,1)}.$ 
		Since the leading terms in \eqref{asym-5} are independent of $\phi$, all
		$\phi$–dependence is contained in the remainder. Differentiating the remainder
		preserves its order:
		\begin{equation}\label{phi-derivative}
			\partial_\phi^k u^{(1)}_1(r,\phi,t)
			= \mathcal{O}\!\left(r^{\delta_1}\right),
			\qquad k=1,2,
		\end{equation}
		because in strip variables $(\eta_1,\eta_2)$ the decaying part satisfies
		$\partial_{\eta_2}^k \widetilde{R}(\eta,t)=\mathcal{O}(e^{\delta_1\eta_1})$, and
		the change of variables $\eta_1=\log r$, $\eta_2=\phi-\pi+\theta_1$ preserves
		this rate.
		
		Since for small $r$ the boundary condition $\partial_\nu u^{(1)}_1=0$ holds on
		$\mathcal{I}_1$, and the outward normal to $\mathcal{I}_1$ in polar coordinates
		coincides with $\mathbf e_\phi$, we therefore have
		$
		\partial_\phi u^{(1)}_1(r,\theta_1,t)=0.
		$
		For $x\in \Gamma_\varepsilon^{(1,1)}$ and such $r$ we have $\phi(x)=\theta_1+\beta_\varepsilon$, where
		$\beta_\varepsilon=\mathcal{O}(\varepsilon^{1/2})$. Hence, by Taylor expansion and \eqref{phi-derivative},
		\begin{equation}\label{phi-derivative+2}
			\partial_\phi u^{(1)}_1(r_\varepsilon,\phi(x),t)
			= \beta_\varepsilon\,\partial^2_{\phi} u^{(1)}_1(r_\varepsilon,\phi^*,t)
			= \mathcal{O}\!\left(\varepsilon^{1/2} r_\varepsilon^{\delta_1}\right),
		\end{equation}
		for some $\phi^*$ between $\theta_1$ and $\phi(x)$. 
		The unit normal to $\Gamma_\varepsilon^{(1,1)}$ is
		\[
		\boldsymbol{\nu}_\varepsilon
		= \sin\beta_\varepsilon\,\mathbf e_r(\phi)
		+ \cos\beta_\varepsilon\,\mathbf e_\phi(\phi),
		\]
		with $\sin\beta_\varepsilon=\mathcal{O}(\varepsilon^{1/2})$ and
		$\cos\beta_\varepsilon=1+\mathcal{O}(\varepsilon)$.
		Thus
		\[
		\partial_{\boldsymbol{\nu}_\varepsilon} u^{(1)}_1
		= \partial_r u^{(1)}_1\,\sin\beta_\varepsilon
		+ \frac{1}{r_\varepsilon}\,\partial_\phi u^{(1)}_1\,\cos\beta_\varepsilon.
		\]
		
		Since $r_\varepsilon=\mathcal{O}(\varepsilon^{1/2})$, the estimates
		\eqref{asym-deriv-r} and \eqref{phi-derivative+2} yield
		\[
		\partial_r u^{(1)}_1\,\sin\beta_\varepsilon
		= C^{(1)}_1(t) + \mathcal{O}\!\left(\varepsilon^{\delta_1/2}\right),
		\qquad
		\frac{1}{r_\varepsilon}\,\partial_\phi u^{(1)}_1
		= \mathcal{O}\!\left(\varepsilon^{\delta_1/2}\right),
		\]
		which proves \eqref{asym-normal-line-final}.
	\end{proof}

\end{document}
